\section{Evolve: Reflective Optimization with an Exact Loss}\label{sec:evolve}

Evolve is the optimizer of milestone M6, specified in \doc{evolve.md} and
\adr{021} to \adr{025}. It searches over plans, physical designs, and whole
pipelines against a recorded workload, with exact replay as its loss and
reconciliation as a hard constraint.

\subsection{Problem}\label{sec:problem}

A pipeline is a DAG of steps. Each step has a logical plan (a DAG of
operators), reads from a physical design (the indexes, layout, and numeric
knobs that exist), and runs many times with different inputs, such as a trigger
batch of keys or a date window. Let $W = \{r_1, \dots, r_n\}$ be the recorded
runs, each with its key set, parameters, and executed plan. An artifact $a$ in
the search space $A$ is a physical design, a plan, a pipeline, or a combination.
Let $c(a, r_i) \in \mathbb{R}^{m}$ be its exact cost vector on $r_i$, the I/O
quantities of Equation~\eqref{eq:cost} extended with measured latency when an
executor is attached; let $m(a)$ be the design's maintenance cost (index bytes
and build time); and let $\Eq(a, a_0)$ be the equivalence predicate of
Section~\ref{sec:nonlocal}. The problem is multi-objective and constrained
(\doc{evolve.md} \S1):
\begin{equation}\label{eq:problem}
  \min_{a \in A}\ \big\{ c(a, r_1), \dots, c(a, r_n),\ m(a) \big\}
  \quad \text{subject to} \quad \Eq(a, a_0).
\end{equation}
No scalar objective is used. Selection is Pareto over runs and objectives
(Section~\ref{sec:pareto}), so a candidate that wins at the median batch and
loses at the p95 batch cannot eliminate one that protects the tail. The loop
is borrowed from text-space optimization~\cite{gepa,skillopt,textgrad} and the
configuration search from physical design~\cite{autoadmin1997}; what is new is
pairing them with an exact loss and a prover-free constraint (\doc{evolve.md}
\S2).

\subsection{The loop}\label{sec:theloop}
\begin{figure}[t]
  \centering
  % Schematic; no measured quantities. Left: the evolve loop (Algorithm 1).
% Right: candidates on two runs, showing why mean selection drops the tail
% protector B while Pareto selection keeps it. Positions are illustrative
% only; the axes carry no scale.
\begin{tikzpicture}[
    font=\scriptsize,
    stage/.style={draw, rounded corners=1.5pt, fill=black!4, inner sep=1.5pt,
                  minimum height=0.62cm, minimum width=0.98cm, align=center},
    flow/.style={-{Stealth[length=1.4mm]}, semithick},
  ]
  % ---- The loop.
  \node[stage] (sel) at (0.55,2.95) {select\\[-1pt]parent};
  \node[stage] (rep) at (1.75,2.95) {replay};
  \node[stage] (att) at (2.95,2.95) {attribute\\[-1pt]$\Delta$, $I$};
  \node[stage] (pro) at (2.95,1.75) {propose};
  \node[stage] (gat) at (1.75,1.75) {gate\\[-1pt]$\Eq(a, a_0)$};
  \node[stage] (upd) at (0.55,1.75) {Pareto\\[-1pt]update};
  \draw[flow] (sel) -- (rep);
  \draw[flow] (rep) -- (att);
  \draw[flow] (att) -- (pro);
  \draw[flow] (pro) -- (gat);
  \draw[flow] (gat) -- (upd);
  \draw[flow] (upd) -- (sel);
  \draw[flow, densely dotted] (gat.south) -- ++(0,-0.42)
    node[below, inner sep=1pt] {discard, with reason};
  \node[text=black!70] at (1.75,2.35) {population $P \ni a_0$};

  % ---- Two runs: cost on the median run against cost on the p95 run.
  \begin{scope}[shift={(4.3,0.45)}]
    % Region dominated by the frontier {A, B}.
    \fill[black!6] (0.6,3.05) -- (0.6,1.7) -- (1.9,1.7) -- (1.9,0.55)
      -- (4.05,0.55) -- (4.05,3.05) -- cycle;
    \draw[bwA, thick] (0.6,3.05) -- (0.6,1.7) -- (1.9,1.7) -- (1.9,0.55)
      -- (4.05,0.55);
    % Equal mean cost through A: x + y is constant along this line.
    \draw[bwInk, densely dashed] (0,2.3) -- (2.3,0);
    \node[rotate=-45, anchor=north, inner sep=1.5pt, text=bwInk]
      at (1.05,1.25) {equal mean};
    % Axes, without a scale.
    \draw[-{Stealth[length=1.4mm]}, bwAxis] (0,0) -- (4.1,0)
      node[pos=0.5, below, text=bwInk, inner sep=1.5pt] {cost on the median run};
    \draw[-{Stealth[length=1.4mm]}, bwAxis] (0,0) -- (0,3.15)
      node[pos=0.5, above, rotate=90, text=bwInk, inner sep=1.5pt]
      {cost on the p95 run};
    % Candidates.
    \node[circle, fill=bwA, inner sep=1.4pt] (A) at (0.6,1.7) {};
    \node[anchor=south west, inner sep=1.5pt] at (A.north east) {$A$: lowest mean};
    \node[fill=bwB, inner sep=1.7pt] (B) at (1.9,0.55) {};
    \node[anchor=south west, inner sep=1.5pt] at (B.north east)
      {$B$: cheapest p95 run};
    \node[draw=black!60, fill=white, regular polygon, regular polygon sides=3,
          inner sep=1.1pt] (C) at (1.3,2.45) {};
    \node[anchor=west, inner sep=1.5pt] at (C.east) {$C$: dominated by $A$};
    \node[star, star points=5, star point ratio=2.3, fill=bwInk,
          inner sep=0.9pt] (a0) at (3.45,2.75) {};
    \node[anchor=east, inner sep=1.5pt] at (a0.west) {baseline $a_0$};
    % What each rule keeps.
    \node[anchor=east, align=right, inner sep=1pt] at (4.05,1.2)
      {mean keeps $A$ only\\Pareto keeps $A$ and $B$};
  \end{scope}
\end{tikzpicture}%

  \caption{Schematic; no measured quantities. Left: the loop of
    Algorithm~\ref{alg:evolve}. Right: candidates scored on two runs. $A$ has
    the lowest mean cost, so mean selection keeps $A$ and drops $B$; $B$ is
    cheaper on the p95 run, so neither dominates the other and Pareto
    selection keeps both, $B$ as the tail protector. $A$ dominates $C$, which
    either rule drops.}
  \Description{Left, a cycle of small boxes: select parent, replay,
    attribute, propose, gate, and Pareto update, with an exit labeled discard
    at the gate. Right, a plot with unnumbered axes, cost on the median run and
    cost on the p95 run, showing the baseline, candidates A, B, and C, a
    staircase frontier through A and B over a shaded dominated region, and a
    dashed line of equal mean cost through A that passes below B.}
  \label{fig:schematic-pareto}
\end{figure}


Algorithm~\ref{alg:evolve} restates \doc{evolve.md} \S3. The baseline is always
in the population. A child that fails the gate is discarded with its reason
recorded; a child that passes is replayed and offered to the frontier. Every
step writes a ledger record, and a proposer's prompt and completion are stored
verbatim with the model identifier. Figure~\ref{fig:schematic-pareto} draws the
loop beside its selection rule.

\begin{algorithm}[t]
  \caption{Evolve (after \doc{docs/evolve.md} \S3). $W_s$ and $W_h$ are the
    search and held-out runs. Every computed value is also written to the
    ledger.}
  \label{alg:evolve}
  \small
  \begin{algorithmic}[1]
    \Require baseline $a_0$; runs $W_s$, $W_h$; rulebook $\mathcal{B}$;
      budget; seed
    \State $P \gets \{a_0\}$ \Comment{the baseline is always a member}
    \While{budget remains}
      \State $p \gets \textsc{Select}(P)$ \Comment{frontier, crowding tie-break}
      \State $\tau \gets \textsc{Replay}(p, W_s)$ \Comment{per run, per operator}
      \State $(\Delta, I) \gets \textsc{Attribute}(p, \tau)$
        \Comment{top-$k$ counterfactuals}
      \State $a \gets \textsc{Propose}(p, \tau, \Delta, I, \mathcal{B})$
        \Comment{one of three proposers}
      \If{not $\Eq(a, a_0)$}
        \State record the discard and its reason; \textbf{continue}
      \EndIf
      \State $C \gets \textsc{Replay}(a, W_s)$
      \State $P \gets \textsc{ParetoUpdate}(P, a, C)$
    \EndWhile
    \State $\mathcal{B} \gets \textsc{Distill}(\text{ledger})$, each rule
      validated on $W_h$
    \State \Return frontier of $P$, receipts, $\mathcal{B}$
  \end{algorithmic}
\end{algorithm}

\subsection{The backward pass: counterfactual attribution}\label{sec:attribution}

\begin{figure}[t]
  \centering
  % Schematic; no measured quantities. The backward pass on a small plan:
% every attributed node carries its counterfactual difference Delta, and one
% pair of choices (u, v) carries a nonzero interaction I_uv.
\begin{tikzpicture}[
    font=\scriptsize,
    op/.style={draw, rounded corners=1.5pt, fill=black!4, inner sep=1.5pt,
               minimum height=0.66cm, minimum width=1.18cm, align=center},
    design/.style={op, fill=white, dashed},
    input/.style={draw=black!55, rounded corners=1.5pt, inner sep=1.5pt,
                  minimum height=0.5cm, align=center, text=black!80},
    flow/.style={-{Stealth[length=1.4mm]}, semithick},
  ]
  % Operators, forward.
  \node[op] (scan) at (0.62,1.55) {scan\\[-1pt]$\Delta_{\mathrm{scan}}$};
  \node[op] (win)  at (2.32,1.55) {window\\[-1pt]$\Delta_{\mathrm{win}}$};
  \node[op, draw=bwA, very thick] (sj) at (4.02,1.55) {semi-join $v$\\[-1pt]$\Delta_v$};
  \node[op] (agg)  at (5.72,1.55) {aggregate\\[-1pt]$\Delta_{\mathrm{agg}}$};
  \node[op, fill=white] (out) at (7.42,1.55) {output\\[-1pt]$c(a)$};
  \draw[flow] (scan) -- (win);
  \draw[flow] (win) -- (sj);
  \draw[flow] (sj) -- (agg);
  \draw[flow] (agg) -- (out);
  % Inputs: a design choice and the run's trigger keys.
  \node[design, draw=bwA, very thick] (key) at (0.62,2.85) {key spec $u$\\[-1pt]$\Delta_u$};
  \node[input] (trig) at (5.95,2.85) {trigger keys $K_r$};
  \draw[flow] (key) -- (scan);
  \draw[flow] (trig.south) -- (sj.north east);
  % The interaction between u and v.
  \draw[bwA, densely dashed, thick] (key.east) to[out=0, in=115]
    (sj.north west);
  \node[anchor=south west, inner sep=1.5pt, text=bwInk] at (1.35,2.95)
    {$I_{uv} \neq 0$: propose $u$ and $v$ together};
  % The backward direction: one exact replay per node.
  \draw[-{Stealth[length=1.6mm]}, black!55, densely dashed]
    (7.42,0.95) -- (0.62,0.95);
  \node[anchor=north, inner sep=1.5pt, align=center] at (4.02,0.92)
    {backward: per node, one exact replay of $a \ominus v$ ($v$ replaced by its
     best local alternative)\\$\Delta_v = c(a) - c(a \ominus v)$};
\end{tikzpicture}%

  \caption{Schematic; no measured quantities. The backward pass on a small
    plan. Solid arrows carry data forward. Each attributed node $v$ carries
    $\Delta_v = c(a) - c(a \ominus v)$, one exact replay of $a$ with $v$
    replaced by its best local alternative. The pair $(u, v)$, a key spec and
    the semi-join, has $I_{uv} \neq 0$: a composite key that pays off only
    once the semi-join moves, so the two changes are proposed together as a
    patch set.}
  \Description{A plan drawn left to right: a key spec feeding a scan, then a
    window, a semi-join that also takes the trigger keys, an aggregate, and
    the output. Each node carries a label Delta sub v, and the output carries
    the cost c of a. A dashed arc between the key spec and the semi-join is
    labeled I sub uv not equal to zero, propose u and v together. A dashed
    arrow runs backward under the plan: one exact replay per node.}
  \label{fig:schematic-backward}
\end{figure}

A gradient says how much each parameter contributed to the loss. For a plan,
the analogue is how much each operator, and each design choice that fixed the
operator's input, contributed to the bytes read. Because replay is exact and
cheap, this is computed by counterfactual rather than approximated
(\adr{024}; Figure~\ref{fig:schematic-backward}). For an artifact $a$ and a
plan node $v$,
\begin{equation}\label{eq:attr}
  \Delta_v = c(a) - c(a \ominus v),
\end{equation}
where $a \ominus v$ replaces $v$ by its best known local alternative (the rule
catalog identity that applies there) or ablates it when none exists. For a pair
of choices $(u, v)$, with $a_u$ denoting $a$ after changing $u$,
\begin{equation}\label{eq:interaction}
  I_{uv} = c(a_{uv}) - c(a_u) - c(a_v) + c(a).
\end{equation}
Both are evaluated per run. A nonzero $I_{uv}$ means the two choices do not
compose additively, the signal that they must be proposed together as a patch
set. Attribution costs one replay per node and the interaction matrix one per
pair. It is capped at the top $k$ nodes by cost rather than sampled, because
sampling would reintroduce estimation, and Shapley values over all subsets were
rejected as exponential (\adr{024}). The design budget for attributing a
20-node plan is one second (\doc{evolve.md} \S12).

\subsection{Three proposers, and what ``gradient'' means}\label{sec:proposers}

The artifacts are discrete: text, DAGs, and sets. There is no differentiable
path from a plan to its cost, so calling the loop gradient descent without
qualification would be inaccurate (\adr{022}). Evolve uses three substitutes
with one scorer and one gate.

\paragraph{Textual gradients for structure.} For operator order, join strategy,
keys, and layouts, the proposer is a language model given the parent, its
per-run trace, its attribution, and the rulebook, and asked for a critique and
one mutation. The critique is a textual gradient in the sense of
ProTeGi~\cite{protegi} and TextGrad~\cite{textgrad}, and the mutation is the
update step. The difference is that the critique is grounded in $\Delta_v$: the
model is told which operator cost what instead of inferring it. Mutations are
plan edits that name an identity when one applies (for example,
\code{semijoin\us before\us window}), so the gate can use a static proof for that
edit and reconciliation for the rest.

\paragraph{Exact numeric descent for knobs.} Numeric knobs include bloom
false-positive rates, key-count gates (when to broadcast, push an \code{IN}
list, or fall back to a shuffle), lookback windows, chunk and segment sizes, and
partition counts. Cost is a deterministic function of them, so the central
difference $\big(c(\theta + h e_j) - c(\theta - h e_j)\big)/2h$ has no
sampling noise, only discretization error, and coordinate descent, pattern
search, and integer line search apply without variance reduction. This is the
one part of evolve that is literally gradient-like, and it is limited to knobs.
Following Section~\ref{sec:exact}, knobs that change what is read from existing
files (rates, gates, lookbacks) are replayed exactly at every probe, while knobs
that change the files (partition counts, and chunk sizes for text sources) are
exact only once each probed point is materialized. Because a proposer's output
is only a candidate, the scalar that descent follows, such as bytes on one
template's p95 run, need not match the selection rule: every point descent
visits is scored on the full objective matrix and gated like any other.

\paragraph{Rule-driven enumeration and deferred surrogates.} The rule catalog's
detectors produce deterministic candidates, such as the composite key a fanout
measurement recommends or a bound on an unbounded lookback, so the model is
never asked to rediscover what measurement already says
(\doc{rule-catalog.md}). A differentiable surrogate trained on replay data would
be an estimator; it is admissible only as a proposer, with replay as the
arbiter, and is deferred (\adr{022}).

\subsection{Sets and pipelines, not calls}\label{sec:sets}

\begin{figure}[t]
  \centering
  % Schematic; no measured quantities. Whole-pipeline equivalence: the
% candidate pushes step five's semi-join on the trigger keys into step two.
% Intermediate outputs differ; the certificate compares only the pipeline
% outputs, over the touched keys (Section 5.4).
\begin{tikzpicture}[
    font=\scriptsize,
    st/.style={draw, rounded corners=1.5pt, fill=black!4, inner sep=1pt,
               minimum height=0.62cm, minimum width=0.6cm, align=center},
    wide/.style={st, minimum width=1.15cm},
    moved/.style={wide, draw=bwA, very thick},
    outn/.style={draw, rounded corners=4pt, fill=white, inner sep=1pt,
                 minimum height=0.62cm, minimum width=0.62cm},
    flow/.style={-{Stealth[length=1.3mm]}, semithick},
  ]
  % Baseline a_0, above.
  \node[anchor=east, inner sep=1pt] at (0.4,2.95) {$a_0$};
  \node[st]    (top1) at (0.75,2.95) {1};
  \node[wide]  (top2) at (1.95,2.95) {2: latest,\\[-1pt]all keys};
  \node[st]    (top3) at (3.10,2.95) {3};
  \node[st]    (top4) at (3.95,2.95) {4};
  \node[wide]  (top5) at (5.10,2.95) {5: semi-join\\[-1pt]on $K_r$};
  \node[outn]  (topO) at (6.25,2.95) {$O_0$};
  % Candidate a, below.
  \node[anchor=east, inner sep=1pt] at (0.4,1.25) {$a$};
  \node[st]    (bot1) at (0.75,1.25) {1};
  \node[moved] (bot2) at (1.95,1.25) {2: latest,\\[-1pt]$K_r$ only};
  \node[st]    (bot3) at (3.10,1.25) {3};
  \node[st]    (bot4) at (3.95,1.25) {4};
  \node[wide]  (bot5) at (5.10,1.25) {5};
  \node[outn]  (botO) at (6.25,1.25) {$O$};
  \foreach \r in {top, bot} {
    \draw[flow] (\r1) -- (\r2);
    \draw[flow] (\r2) -- (\r3);
    \draw[flow] (\r3) -- (\r4);
    \draw[flow] (\r4) -- (\r5);
    \draw[flow] (\r5) -- (\r O);
  }
  % Intermediate outputs differ: a per-step check rejects the candidate.
  \foreach \x in {2.66,3.52,4.39} {
    \node[text=bwB, font=\footnotesize] at (\x,2.1) {$\neq$};
  }
  \node[text=bwB, anchor=east, inner sep=1pt] at (2.42,2.1) {per step: reject};
  % Pipeline outputs agree over the touched keys: the certificate passes.
  \node[text=bwA, font=\footnotesize] at (6.25,2.1) {$=$};
  \node[draw=bwA, thick, rounded corners=1.5pt, fill=white, inner sep=1.5pt,
        align=center, anchor=west] at (6.72,2.1)
    {output\\[-1pt]certificate over\\[-1pt]touched keys:\\[-1pt]pass};
  % The move.
  \draw[flow, bwA, thick] (bot5.south) to[out=-155, in=-25]
    node[pos=0.5, below, inner sep=1.5pt, text=bwInk]
      {semi-join on the trigger keys $K_r$ pushed from step 5 into step 2}
    (bot2.south);
\end{tikzpicture}%

  \caption{Schematic; no measured quantities. The candidate moves step five's
    semi-join on the trigger keys into step two, as a temporal \emph{latest}
    lookup (Section~\ref{sec:nonlocal}). The outputs of steps two to four
    differ, so a per-step check rejects it; the certificate compares only the
    pipeline outputs, over the touched keys.}
  \Description{Two rows of five pipeline steps, the baseline above and the
    candidate below, joined by arrows from step to step. A curved arrow moves
    the semi-join from step five to step two. Not-equal signs mark the
    outputs of steps two to four; an equals sign and a certificate mark the
    two pipeline outputs, compared over the touched keys.}
  \label{fig:schematic-pipeline}
\end{figure}

The unit of optimization is never a single call (\adr{023}). Plans are mutated
as DAGs and scored whole, so interactions inside a plan are captured by
construction. Physical design is a set-valued configuration of key specs, a
layout, and knob values, scored by Equation~\eqref{eq:design} against every
run, with prefix-seek subsumption modeled (Section~\ref{sec:counterfactual})
and $m(a)$ kept as a separate objective. This is the classical index-selection
problem~\cite{autoadmin1997,bruno2005relax}, whose weak point was the estimate.
Search is greedy construction followed by local search with exact scoring,
seeded by rule-driven candidates and mutated by the model, and interaction terms
between key specs decide which specs are proposed together. Pipelines are DAGs
of steps with lineage; cross-step mutations are allowed and gated at the
pipeline output (Section~\ref{sec:nonlocal}; Figure~\ref{fig:schematic-pipeline}),
which is what distinguishes optimizing a pipeline from optimizing each of its
queries.

\subsection{Pareto selection over runs}\label{sec:pareto}

Scalarizing cost, for example as mean bytes over runs, hides the tail. For a
candidate $a$ let $C(a)$ be its matrix of runs by objectives, together with
$m(a)$. Candidate $a$ dominates $b$ if it is at least as good in every cell and
strictly better in one, and the frontier is the set of non-dominated members
(\doc{evolve.md} \S6). In particular, a candidate that is strictly best on any
single run survives, which is the rule stated in \adr{021}, and a candidate
that wins on the mean cannot eliminate one that protects the tail
(Figure~\ref{fig:schematic-pareto}). Crowding distance,
as in NSGA-II~\cite{nsga2}, breaks ties when choosing parents so the frontier
stays spread. Runs are grouped by template (same plan shape, different keys) so
that the frontier is not dominated by run count, and $m(a)$ is on the frontier,
so an index that saves bytes read at a large storage cost is visible rather than
hidden. Because the baseline is always in the population, the frontier never
regresses below it.

\subsection{Determinism and the ledger}\label{sec:determinism}

A search is pinned to a manifest sequence, a workload hash, a rule catalog
version, proposer identifiers, and a seed (\doc{evolve.md} \S8). Every
candidate, score, certificate, attribution, interaction matrix, frontier
snapshot, model prompt and completion, and rule is a ledger record, and the
planned contract adds a record kind for each (\code{Candidate},
\code{Attribution}, \code{Interaction}, \code{Frontier}, \code{ProposerCall},
\code{Rule}, \code{Rulebook}; \doc{evolve.md} \S10). Model outputs are not assumed to be
deterministic. Replaying a search from its ledger reuses the stored completions
and reproduces the frontier byte for byte, which extends invariant~3; a fresh
search is a new ledger. Budgets for replays, model calls, and wall-clock time
are explicit, and exhausting one is a recorded terminal state, not an error.

\subsection{The proposer in the threat model}\label{sec:proposerthreat}

Per \doc{evolve.md} \S9, a proposer may be wrong, adversarial, or compromised.
It can waste budget, propose non-equivalent artifacts, which the gate discards, and
write misleading rule text, which held-out validation rejects. It cannot score,
pass the gate, apply, or promote a rule, and a surrogate would inherit the same
limits. Because the loop invokes a language model as one of its proposers,
\adr{009}, as amended, states the boundary precisely: the engine
never consults a model for a score, a verification outcome, or an apply
decision. A proposer, model or otherwise, is invoked only to produce a Proposal
through the same path an external agent uses; its prompt and completion are
ledger records, and nothing it returns is trusted. The engine may host an
untrusted proposer, but the direction of trust is unchanged.

\subsection{The rulebook and its validation}\label{sec:rulebook}

The frontier answers for one workload; the rulebook is what transfers
(\doc{evolve.md} \S7). After a search, winning mutations are clustered by
identity and pattern, and for each cluster the model writes a rule: a detector
over evidence or plans, an action, an invariant, and the evidence hashes and
measured gains that justify it. No rule is accepted on the model's word
(\adr{025}). A rule is promoted only if applying it to held-out runs, excluded
from the search, reproduces a gain under exact replay and passes the gate.
Rules carry a confidence derived from that validation, hash links to the
evidence and receipts that produced them, and their certificates' coverage, and
a rule whose evidence has drifted is flagged stale. This is the validation gate
of SkillOpt~\cite{skillopt} applied to rules rather than skills. The rulebook is
a text artifact that agents read on the next \code{fix}, a learned analogue of
a hand-written optimization catalog, and it is the subject of the transfer
experiment in Section~\ref{sec:eval}.
